\subsection{Zufallsvariablen}
\begin{emphbox}
    Eine Zufallsvariable bildet jedes Ergebnis in einer Ergebnismenge $\Omega$ auf Ergebnisse in $\Omega'$ ab.
    Für reelle Zufallsvariablen gilt $\Omega'=\mathbb R^n$.
    \begin{equation*}
        X : \Omega \rightarrow \Omega'
    \end{equation*}
\end{emphbox}

\begin{sectionbox}
    \subsubsection{Messraum}
    Zufallsvariablen können aus einem abstrakten Raum $\Omega$ ein Ergebnis $\omega\in\Omega$ auf einen Wert $X(\omega)\in\mathbb R$ abbilden.
    Der Raum $X(\Omega)$ ist der Messraum, wobei der abstrakte Raum viele Dimensionen hat und schwierig zu vermessen wäre (z.B. Kreidewurf, Ergebnis, dass die Kreide genau dort landet: Wurfwinkel, Geschwindigkeit, Raumtemperatur, ...)
\end{sectionbox}